\chapter{Phonon-qubit coupling via piezoelectricity}
\label{ch:phonon-qubit-piezo}

The complementary recoil-channel readout, parallel to the
phonon-quasiparticle conversion of \cref{ch:qp-detection}, is to
read out individual phonons in the substrate directly --- before
they break Cooper pairs or thermalise --- by coupling them to a
superconducting qubit through the piezoelectric response of the
substrate material. The same physics is the workhorse of
\emph{quantum acoustics}: a body of work in which superconducting
transmons are coherently entangled with single phonons of GHz
acoustic resonators, with the phonon playing the role that a
single microwave photon plays in cavity QED. The chapter starts
with the elastic-wave background needed to enumerate phonon modes
in the relevant geometries, surveys the experimental quantum
acoustics literature, and then maps the same coupling onto the
recoil pipeline of \cref{ch:dm-search} as a single-phonon
dark-matter sensor.

% --------------------------------------------------------------
\section{Phonons and strain in solids}
\label{sec:phonon-modes}

\begin{phenomenon}
A solid body responds to small deformations as a continuum elastic
medium. Its normal modes of vibration are quantised as phonons,
each carrying energy $\hbar\omega$ at angular frequency $\omega$.
The mode spectrum is fixed by the body's geometry, by its
elastic constants, and by the boundary conditions on its
surfaces.
\end{phenomenon}

\begin{statement}[Elastic wave equation and mode spectrum]
\label{stmt:elastic-modes}
For a solid of mass density $\rho$ and elastic-stiffness tensor
$c_{ijkl}$, the displacement field $\bm u(\bm r,t)$ satisfies the
elastic wave equation~\cite{tinkham2004}
\begin{equation}
\label{eq:elastic-wave}
\rho\,\frac{\partial^2 u_i}{\partial t^2}
\;=\; \frac{\partial}{\partial x_j}
\left[c_{ijkl}\,\frac{\partial u_k}{\partial x_l}\right],
\end{equation}
with appropriate boundary conditions: free surfaces require
$\sigma_{ij}\hat n_j = c_{ijkl}\partial_l u_k\,\hat n_j = 0$ (no
applied stress), clamped surfaces require $\bm u = 0$. Phonon
modes are the time-harmonic solutions
$\bm u(\bm r,t) = \bm u_n(\bm r) e^{-i\omega_n t}$ with mode
profile $\bm u_n$ satisfying the eigenvalue equation
$\rho \omega_n^2 u_{n,i} = -\partial_j(c_{ijkl}\partial_l u_{n,k})$.
\end{statement}

\paragraph{Mode families.}
Three families of modes recur in the structures used for cQED
quantum acoustics and for cryogenic dark-matter detection:
\begin{itemize}
  \item \emph{Bulk longitudinal and transverse acoustic modes.}
        In an unbounded crystal, plane-wave solutions of
        \cref{eq:elastic-wave} give the acoustic-phonon
        dispersion $\omega(\bm k)\approx c_s |\bm k|$ at low
        $|\bm k|$, with longitudinal and transverse polarisation
        propagating at sound speeds $c_L,\,c_T$. In silicon,
        $c_L\approx 8.4\,\text{km/s}$ and
        $c_T\approx 5.8\,\text{km/s}$; in sapphire,
        $c_L\approx 11\,\text{km/s}$.
  \item \emph{Surface acoustic waves (SAWs, Rayleigh modes).}
        On a free surface, an additional family of modes lives
        near the surface and decays exponentially into the bulk
        on a length scale of one wavelength. SAWs propagate
        along the surface at a speed slightly below $c_T$ and
        have energy concentrated within a few $\lambda$ of the
        surface, which makes them efficient at coupling to
        surface-mounted electrodes.
  \item \emph{Confined modes of a finite body.}
        For a body of finite extent (slab, beam, dome), the
        spectrum becomes discrete with mode spacing
        $\Delta\omega\sim c_s/L$. Thin-film thickness modes of an
        AlN or GaAs film of thickness $d\sim\mu\text{m}$ live at
        $\omega/2\pi\sim c_s/(2d)\sim$ GHz, exactly the band of
        a transmon qubit, which is why HBAR (high-overtone bulk
        acoustic resonator) and SAW devices are the standard
        platforms for quantum acoustics.
\end{itemize}

\paragraph{Density of states.}
For a bulk solid the low-energy density of states is the Debye
form $\rho_{\rm ph}(\omega)\propto\omega^2$ up to a cutoff at the
Debye frequency $\omega_D=k_B T_D/\hbar$. Optical phonon branches
in polar crystals (e.g.\ the LO-TO splitting in GaAs) sit at
finite energy near the zone centre and are particularly relevant
for sub-GeV DM scattering~\cite{knapen2018}, because their dipole
moments couple efficiently to the electromagnetic response of
the medium and because their energies of $\sim$ tens of meV are
matched to the DM kinematic threshold.

% --------------------------------------------------------------
\section{Quantum acoustics with superconducting qubits}
\label{sec:quantum-acoustics}

\begin{phenomenon}
A single phonon of an acoustic resonator on a piezoelectric
substrate carries an electric polarisation. Coupled to the charge
degree of freedom of a transmon qubit, this polarisation acts as
an effective electric dipole interaction, with the phonon playing
the role that a single microwave photon plays in cavity QED. In
the strong-coupling regime, the qubit and the phonon undergo
coherent vacuum-Rabi oscillations and form the acoustic analogue
of dressed states.
\end{phenomenon}

\begin{statement}[Phonon-qubit Jaynes--Cummings coupling]
\label{stmt:phonon-qubit-jc}
On a piezoelectric substrate with electromechanical coupling
$\hat e_{14}$, the linear coupling between a single confined
phonon mode of frequency $\omega_{\rm ph}$ and a transmon at
frequency $\omega_q$ takes, in the rotating-wave approximation
(\cref{sec:rwa}), the Jaynes--Cummings form
\begin{equation}
\label{eq:phonon-qubit-jc}
\op{H}_{\rm int} \;=\; \hbar g\,
\bigl(\hat b\,\hat\sigma_+ \,+\, \hat b^\dagger\,\hat\sigma_-\bigr),
\end{equation}
with $\hat b,\hat b^\dagger$ the phonon mode operators and
$g$ the single-phonon coupling strength
\begin{equation}
\label{eq:phonon-coupling}
g \;\propto\; e_{14}\,d\,\sqrt{\frac{\hbar\,\omega_{\rm ph}}{2\,M\,\Omega_{\rm ph}}},
\end{equation}
where $d$ is the spatial-overlap length between the phonon mode
and the qubit electrode and $M$ is the effective phonon-mode
mass. Typical experimental values are
$g/2\pi\sim 0.1\text{--}10\,\text{MHz}$, well within the
strong-coupling regime $g\gg\kappa_{\rm ph},\,1/T_2$.
\end{statement}

\paragraph{Two device families: SAW and HBAR.}
Two architectures dominate the quantum-acoustics literature.
\begin{itemize}
  \item \emph{Surface acoustic-wave devices.}
        An interdigital transducer (IDT) on a piezoelectric
        substrate (LiNbO$_3$, GaAs, AlN) launches and receives
        Rayleigh waves; coupling to a transmon comes from the IDT
        electrodes capacitively driving the qubit. Gustafsson
        et~al.~\cite{gustafsson2014} demonstrated the first
        artificial-atom-SAW system in the strong-coupling regime.
        Manenti et~al.~\cite{manenti2016} demonstrated SAW
        resonators with quality factors approaching $10^5$ in the
        single-phonon limit. Satzinger et~al.~\cite{satzinger2018}
        achieved deterministic single-phonon control on a SAW
        resonator using a flux-tunable transmon.
  \item \emph{High-overtone bulk acoustic resonators.}
        A piezoelectric thin film (typically AlN, $\sim\mu$m
        thick) is sandwiched between two electrodes; the film
        supports thickness-mode bulk-acoustic standing waves at
        $\omega/2\pi\sim$ GHz with extraordinarily high quality
        factors ($Q\gtrsim 10^{6}$ at low temperature) because the
        confining surfaces are atomically flat. Chu
        et~al.~\cite{chu2017} demonstrated the first
        Jaynes--Cummings coupling between a transmon and an HBAR
        single-phonon mode; Chu~et~al.~\cite{chu2018} subsequently
        showed deterministic preparation and tomography of
        multi-phonon Fock states up to $|7\rangle$. The HBAR is
        essentially the acoustic analogue of a 3D microwave
        cavity for cQED, but with $10^5$ times shorter wavelength
        (a single phonon in a 1\,GHz HBAR fits in 5\,$\mu$m of
        sapphire), enabling extreme miniaturisation.
\end{itemize}
The original demonstration that a discrete mechanical mode could
be cooled to the quantum ground state and operated as a single
quantum object was the dilatational-mode dilation experiment of
O'Connell et~al.~\cite{oconnell2010}, which set the template that
the present cQED-coupled quantum-acoustics platforms inherited.

% --------------------------------------------------------------
\section{Dark-matter detection via single-phonon readout}
\label{sec:phonon-dm}

\begin{phenomenon}
A dark-matter particle that scatters in a piezoelectric crystal
deposits energy as a population of athermal phonons. Each phonon
that survives long enough to reach a coupled qubit can be detected
as a $\ket{0}\to\ket{1}$ transition --- exactly as a single
microwave photon is detected in cavity QED, but with phonon
masses of order an MeV per quantum rather than 0 for a photon. The
recoil pipeline of \cref{sec:fermion-dm} combined with the
single-phonon readout of \cref{sec:quantum-acoustics} therefore
gives a sub-eV-threshold direct-detection scheme for sub-GeV DM.
\end{phenomenon}

\begin{statement}[Single-phonon DM detection]
\label{stmt:phonon-dm}
For a target of total mass $\mathcal M$ on a piezoelectric
substrate that supports a confined phonon mode of frequency
$\omega_{\rm ph}$ coupled to a transmon at the same frequency
with single-phonon coupling $g$, an event of recoil energy
$E_R=\hbar\omega_{\rm ph}$ on resonance has a phonon-creation
probability
\begin{equation}
\label{eq:phonon-eta}
\eta_{\rm ph}(E_R) \;\sim\; \eta_{\rm coll}(E_R)\;\frac{|\langle 1|\hat b^\dagger|0\rangle|^2}{\bigl(E_R-\hbar\omega_{\rm ph}\bigr)^2/(\hbar\,\Gamma_{\rm ph})^2 + 1},
\end{equation}
which peaks for resonant deposition. Off-resonance deposition
produces phonons at the corresponding bath energies, which
thermalise and diffuse out of the readout band; the on-resonance
sensitivity is therefore narrow-band, but optimised for whichever
phonon mode the qubit is tuned to.
\end{statement}

\paragraph{Mass reach.}
For a meV-scale phonon mode and a kg-scale piezoelectric target
(silicon-on-AlN composite, sapphire-on-GaAs heterostructure, or
similar), the kinematic threshold of \cref{eq:E_R-max} for
$E_R=\hbar\omega_{\rm ph}$ gives
$m_\chi^{\rm min}\sim\sqrt{M_T E_R}/v_{\rm esc}$. With
$M_T = 28\,m_p$ (Si nucleus) and
$E_R=10\,\text{meV}\sim 10^{-11}\,\text{GeV}$,
$m_\chi^{\rm min}\sim 6\,\text{MeV}$. Linehan
et~al.~\cite{linehan2024} present the detailed analysis for a
realistic device including substrate phonon-mode density and
qubit-readout efficiency, and project a reach of order
$\sigma_n\lesssim 10^{-38}\,\text{cm}^2$ at
$m_\chi\sim 1$--$10\,\text{MeV}$ for a kg-day exposure. These
numbers compare favourably with the phonon-quasiparticle
detectors of \cref{ch:qp-detection} in the same mass band, and
together the two approaches define the sub-GeV reach plot of the
contemporary cQED dark-matter programme.

\paragraph{Light-DM advantages of piezoelectric coupling.}
Two features make the piezoelectric approach particularly
attractive at low DM mass:
\begin{itemize}
  \item \emph{Single-phonon readout, not energy resolution.}
        A transmon coupled to an acoustic mode gives a
        \emph{binary} signal --- the phonon mode is occupied or
        not --- so the threshold is set by the qubit dark-count
        rate, typically $\Gamma_{\rm dc}\sim 1\,\text{mHz}$, and
        not by an analog energy-resolution floor as in a TES.
  \item \emph{Optical-phonon coupling.}
        For polar crystals, the optical-phonon dipole moment
        couples directly to the DM-electron and DM-photon
        interaction channels~\cite{knapen2018}, opening detection
        modes for dark photons, axion-like dark matter, and
        light-mediator DM-electron scattering that are not
        accessible to non-polar absorbers.
\end{itemize}

\paragraph{Limitations and outlook.}
The narrow-band character of \cref{eq:phonon-eta} is the main
constraint: a single qubit-phonon system covers a narrow
$\Delta\omega/\omega\sim 1/Q$ slice of the recoil-energy spectrum.
Realistic experiments either tile multiple qubit-phonon systems at
different frequencies to broaden the response or scan the qubit
frequency in time. Both strategies are demonstrated in present
quantum-information hardware (the same multiplexing solutions
used for KIDs), so the scaling path to a kg-day-class device is
identical to the scaling path of a multi-thousand-qubit processor.
This is the deepest cross-talk between the dark-matter and
fault-tolerant-computing programmes that drives the
chapter-6 / chapter-8 connection.
